\documentclass[a4paper,12pt]{article}
\usepackage[T2A]{fontenc}
\usepackage[utf8]{inputenc}
\usepackage[russian]{babel}
\usepackage{amsmath,amssymb,mathtools}
\usepackage{PTSans}
\renewcommand{\familydefault}{\sfdefault}
\usepackage{icomma}
\usepackage[a4paper,left=23mm,right=23mm,top=22mm,bottom=23mm]{geometry}
\usepackage{microtype}
\usepackage{tikz}
\usetikzlibrary{calc,arrows.meta,positioning,shapes.geometric}
\usepackage{tkz-euclide}
\usepackage{tabularx,booktabs}
\usepackage{enumitem}
\usepackage{fancyhdr}
\usepackage{setspace}
\usepackage{xcolor}
\usepackage[unicode,hidelinks]{hyperref}
\definecolor{accent}{HTML}{087B79}
\definecolor{darkaccent}{HTML}{164D55}
\definecolor{textgray}{HTML}{37454A}
\definecolor{lightaccent}{HTML}{D8E8E6}
\color{textgray}
\onehalfspacing
\setlength{\parindent}{0pt}
\setlength{\parskip}{4pt}
\setlist[itemize]{leftmargin=*,itemsep=2pt,topsep=3pt}
\setlist[enumerate]{leftmargin=*,itemsep=7pt,topsep=5pt}
\setcounter{secnumdepth}{0}
\pagestyle{fancy}
\fancyhf{}
\fancyhead[L]{\footnotesize\color{darkaccent}Екатерина Скелли · 17.10.2026}
\fancyhead[R]{\footnotesize\color{darkaccent}Повторение занятия}
\fancyfoot[C]{\footnotesize\thepage}
\renewcommand{\headrulewidth}{0.25pt}
\renewcommand{\footrulewidth}{0pt}
\newcommand{\key}[1]{\textbf{\color{darkaccent}#1}}
\newcommand{\checkitem}[1]{\item[\(\square\)] #1}
\newenvironment{skills}{\begin{itemize}[label=\(\square\),leftmargin=*,itemsep=2pt]}{\end{itemize}}
\newcommand{\sys}[2]{\left(#1\right)_{#2}}
\begin{document}

% ---------- Титульный лист ----------
\thispagestyle{empty}
\begin{center}
\vspace*{30mm}
{\color{accent}\Large МАТЕМАТИКА · ГЕОМЕТРИЯ · СИСТЕМЫ СЧИСЛЕНИЯ}\par
\vspace{19mm}
{\fontsize{32}{38}\selectfont\bfseries\color{darkaccent}Чек-лист}\par
\vspace{10mm}
{\fontsize{20}{27}\selectfont Параллелограмм и прямоугольник}\par
\vspace{3mm}
{\Large Позиционные системы счисления}\par
\vspace{22mm}
{\large 17 октября 2026 года}\par
\vspace{8mm}
{\large Екатерина Скелли}\par
\vfill
{\small Лёвин Артём Александрович}\par
\vspace{2mm}
{\small эксклюзивно для Екатерины Скелли}\par
\end{center}
\vspace{18mm}
\begin{tikzpicture}[remember picture,overlay]
  \draw[accent!55,line width=1.3pt] ([xshift=18mm,yshift=21mm]current page.south west)
    .. controls ([xshift=69mm,yshift=34mm]current page.south west)
    and ([xshift=-75mm,yshift=5mm]current page.south east)
    .. ([xshift=-18mm,yshift=24mm]current page.south east);
\end{tikzpicture}
\clearpage

\section*{1. Параллелограмм: определение и свойства}
\key{Параллелограмм} --- четырёхугольник, у которого противоположные стороны попарно параллельны.

\begin{center}
\begin{tikzpicture}[scale=0.93]
\tkzDefPoint(0,0){A}
\tkzDefPoint(5.2,0){B}
\tkzDefPoint(6.5,2.5){C}
\tkzDefPoint(1.3,2.5){D}
\tkzDefMidPoint(A,C)\tkzGetPoint{O}
\tkzDrawPolygon[color=darkaccent,line width=1pt](A,B,C,D)
\tkzDrawSegments[color=accent,line width=.85pt](A,C B,D)
\tkzDrawPoints[size=2.6,color=darkaccent](A,B,C,D,O)
\tkzLabelPoints[below](A,B)
\tkzLabelPoints[above](C,D)
\tkzLabelPoints[below right](O)
\end{tikzpicture}
\end{center}

\begin{skills}
\checkitem{Узнаю свойства сторон: \(AB=CD\), \(BC=AD\).}
\checkitem{Использую свойства углов: \(\angle A=\angle C\), \(\angle B=\angle D\); сумма соседних углов равна \(180^\circ\).}
\checkitem{Нахожу периметр по двум соседним сторонам: \(P=2(AB+BC)\).}
\checkitem{Помню свойство диагоналей: \(AO=OC=\dfrac{AC}{2}\), \(BO=OD=\dfrac{BD}{2}\).}
\checkitem{Понимаю: диагонали произвольного параллелограмма могут быть разной длины.}
\end{skills}

\key{Почему противоположные стороны равны?} Проведём диагональ \(AC\). Из параллельности противоположных сторон получаем \(\angle BAC=\angle DCA\) и \(\angle BCA=\angle DAC\). Сторона \(AC\) общая, поэтому \(\triangle BAC\cong\triangle DCA\) по стороне и двум прилежащим к ней углам. Следовательно, \(AB=CD\) и \(BC=AD\).

\subsection*{Два коротких примера}
\textbf{Пример 1.} \(P=54\) см, \(AB:BC=5:4\). Обозначим \(AB=5x\), \(BC=4x\). Тогда
\[
  2(5x+4x)=54,\qquad x=3,\qquad AB=15\ \text{см},\quad BC=12\ \text{см}.
\]
\textbf{Пример 2.} \(\angle A:\angle B=2:3\). Соседние углы дополняют друг друга до \(180^\circ\):
\[
2x+3x=180^\circ,\quad x=36^\circ,\quad \angle A=72^\circ,\quad\angle B=108^\circ.
\]
\key{Важно.} Отношение \(5:4\) означает \(5x\) и \(4x\), а не длины \(5\) и \(4\).

\clearpage
\section*{2. Как доказывать, что четырёхугольник --- параллелограмм}
\key{Признаки} применяют, когда принадлежность фигуры к параллелограммам ещё предстоит доказать. \key{Свойства} применяют после того, как это уже установлено.

Для четырёхугольника достаточно \emph{любого одного} из трёх признаков:
\begin{enumerate}[label=\textbf{\arabic*.},itemsep=3pt]
\item \key{Одна пара противоположных сторон одновременно равна и параллельна.}
\item \key{Обе пары противоположных сторон равны.}
\item \key{Диагонали точкой пересечения делятся пополам.}
\end{enumerate}

\subsection*{Пример доказательства с серединами}
В параллелограмме \(ABCD\) точки \(M\) и \(N\) --- середины сторон \(AB\) и \(CD\) соответственно. Докажем, что \(AMCN\) --- параллелограмм.

\(AM\parallel CN\), поскольку \(AB\parallel CD\). Кроме того,
\[
 AM=\frac{AB}{2}=\frac{CD}{2}=CN.
\]
Одна пара противоположных сторон четырёхугольника \(AMCN\) равна и параллельна. \key{Значит, \(AMCN\) --- параллелограмм (признак 1).}

\subsection*{Пример доказательства с диагоналями}
Пусть \(O\) --- пересечение диагоналей параллелограмма \(ABCD\), а точки \(E\) и \(F\) расположены на \(AC\) в порядке \(A,E,O,F,C\), причём \(AE=CF\). Докажем, что \(BEDF\) --- параллелограмм.
\[
 AO=OC,\quad AE=CF \quad\Longrightarrow\quad
 EO=AO-AE=OC-CF=OF.
\]
Также \(BO=OD\). Следовательно, диагонали \(BD\) и \(EF\) четырёхугольника \(BEDF\) пересекаются и делятся пополам. \key{Применяем признак 3.}

\subsection*{Середины сторон произвольного четырёхугольника}
Если последовательно соединить середины сторон \emph{любого} четырёхугольника, получится параллелограмм. Для доказательства проведите диагональ исходного четырёхугольника: две противоположные стороны внутренней фигуры окажутся средними линиями двух треугольников, а значит, будут равны половине одной диагонали и параллельны ей.

\clearpage
\section*{3. Биссектрисы и прямоугольник}
\subsection*{Биссектрисы в параллелограмме}
\begin{skills}
\checkitem{Если биссектрисы двух \emph{соседних} внутренних углов пересекаются в точке \(E\), то угол между ними равен \(90^\circ\).}
\checkitem{Если биссектриса угла \(DAB\) пересекает \emph{сторону} \(CD\) в точке \(K\), то \(DK=AD\).}
\end{skills}

\textbf{Почему биссектрисы перпендикулярны?}
\[
\angle AEB=180^\circ-\tfrac12(\angle A+\angle B)
=180^\circ-\tfrac12\cdot180^\circ=90^\circ.
\]
Здесь \(AE\) и \(BE\) --- биссектрисы соседних углов, \(E\) --- их точка пересечения.

\textbf{Почему \(DK=AD\)?} Поскольку \(AB\parallel DK\), имеем \(\angle KAB=\angle AKD\). Биссектриса даёт \(\angle KAB=\angle DAK\). Поэтому \(\angle AKD=\angle DAK\), и треугольник \(ADK\) равнобедренный: \(DK=AD\).

\key{Оговорка.} Такой вывод относится к случаю \(K\in CD\). Биссектриса вообще не обязана делить противоположную сторону пополам; точка \(K\) на стороне \(CD\) существует при \(AB\ge AD\).

\subsection*{Прямоугольник}
\key{Прямоугольник} --- параллелограмм, у которого все углы прямые.

\begin{skills}
\checkitem{Помню: противоположные стороны прямоугольника равны, соседние перпендикулярны; диагонали равны и делятся точкой пересечения пополам.}
\checkitem{Знаю признаки: (1) параллелограмм с одним прямым углом; (2) параллелограмм с равными диагоналями.}
\checkitem{Понимаю необходимость сначала обосновать, что фигура является параллелограммом.}
\end{skills}

\textbf{Доказательство признака 2.} Пусть \(ABCD\) --- параллелограмм и \(AC=BD\). Треугольники \(ABC\) и \(BAD\) равны по трём сторонам: \(AB\) общая, \(BC=AD\), \(AC=BD\). Тогда \(\angle ABC=\angle BAD\). Они соседние углы параллелограмма, их сумма \(180^\circ\), следовательно, каждый равен \(90^\circ\). Значит, \(ABCD\) --- прямоугольник.

\key{Методический ориентир.} В задачах на доказательство сначала ищите подходящие треугольники: проведите диагональ, высоту, медиану или среднюю линию. Любое утверждение о равенстве углов должно опираться на уже установленные условия.

\clearpage
\section*{4. Системы счисления: что обязательно понимать}
\key{Основание системы счисления} определяет допустимые цифры и веса разрядов. Справа налево степени основания начинаются с нуля.

\begin{center}
\renewcommand{\arraystretch}{1.25}
\begin{tabularx}{\linewidth}{@{}lX@{}}
\toprule
Система & Допустимые цифры \\
\midrule
Двоичная, основание 2 & \(0,1\) \\
Восьмеричная, основание 8 & \(0,1,2,3,4,5,6,7\) \\
Десятичная, основание 10 & \(0,1,2,3,4,5,6,7,8,9\) \\
Шестнадцатеричная, основание 16 & \(0,1,\ldots,9,A,B,C,D,E,F\) \\
\bottomrule
\end{tabularx}
\end{center}
\[
 A=10,\qquad B=11,\qquad C=12,\qquad D=13,\qquad E=14,\qquad F=15.
\]
\textbf{Проверка допустимости записи.} \(\sys{76}{8}\) допустимо, а \(\sys{78}{8}\) --- нет: цифры 8 в восьмеричной системе нет. Аналогично \(\sys{1G}{16}\) недопустимо.

\subsection*{Перевод в десятичную систему}
\key{Правило:} умножить каждую цифру на основание в степени, равной её позиции справа (начиная с 0), и сложить.
\[
\begin{aligned}
 \sys{123}{16}&=1\cdot16^2+2\cdot16+3=\sys{291}{10},\\
 \sys{A5}{16}&=10\cdot16+5=\sys{165}{10},\\
 \sys{151}{8}&=1\cdot8^2+5\cdot8+1=\sys{105}{10}.
\end{aligned}
\]
\subsection*{Перевод из десятичной системы}
Делите число с остатком на нужное основание до получения нулевого частного. \key{Читайте остатки снизу вверх.}
\[
\begin{array}{rcl@{\qquad}l}
165&=&8\cdot20+5 & \text{остаток }5,\\
20&=&8\cdot2+4  & \text{остаток }4,\\
2&=&8\cdot0+2   & \text{остаток }2.
\end{array}
\]
Итак, \(\sys{165}{10}=\sys{245}{8}\). Обратная проверка:
\(2\cdot8^2+4\cdot8+5=165\).

\clearpage
\section*{5. Переводы через двоичную систему}
\key{Триады} --- группы по 3 двоичные цифры для перехода между системами с основаниями \(2\) и \(8\). \key{Тетрады} --- группы по 4 двоичные цифры для перехода между системами с основаниями \(2\) и \(16\).

\begin{center}
\renewcommand{\arraystretch}{1.2}
\begin{tabular}{@{}cccccc@{}}
\toprule
16-рич. & Двоичная тетрада & & 16-рич. & Двоичная тетрада \\
\midrule
\(A\) & \(1010\) & & \(D\) & \(1101\) \\
\(B\) & \(1011\) & & \(E\) & \(1110\) \\
\(C\) & \(1100\) & & \(F\) & \(1111\) \\
\bottomrule
\end{tabular}
\end{center}
\textbf{Пример 1.} \(\sys{3D5}{16}\) представим по тетрадам:
\[
\underbrace{0011}_{3}\ \underbrace{1101}_{D}\ \underbrace{0101}_{5}.
\]
В двоичной записи \(2+3+2=7\) единиц.

\textbf{Пример 2.} Из двоичной в восьмеричную:
\[
\sys{001\ 011\ 010}{2}=\sys{132}{8}.
\]
\textbf{Пример 3.} Из шестнадцатеричной в восьмеричную:
\[
\begin{aligned}
\sys{794}{16}&=\sys{0111\ 1001\ 0100}{2}\\
&=\sys{011\ 110\ 010\ 100}{2}=\sys{3624}{8}.
\end{aligned}
\]

\begin{skills}
\checkitem{При разбиении двоичного числа на группы иду \emph{справа налево}.}
\checkitem{При необходимости дополняю нулями \emph{слева}, не меняя значение числа.}
\checkitem{Не путаю двоичную тетраду \(1010_2\) с десятичным числом 1010: это \(A_{16}\).}
\checkitem{При записи ответа из шестнадцатеричной системы использую латинские \(A\)--\(F\).}
\checkitem{Проверяю результат обратным переводом или через десятичную систему.}
\end{skills}

\key{Как выбрать способ?} Для \(2\leftrightarrow8\) используйте триады, для \(2\leftrightarrow16\) --- тетрады. Для \(8\leftrightarrow16\) удобно сначала перейти в двоичную систему, затем заново сгруппировать разряды. Универсальная альтернатива --- перевод через десятичную систему.

\clearpage
\section*{Алгоритм: как выбрать признак параллелограмма}
\vspace{3mm}
\begin{center}
\begin{tikzpicture}[
  node distance=9mm and 8mm,>=Stealth,
  start/.style={ellipse,draw=darkaccent,thick,align=center,text width=42mm,minimum height=10mm,fill=lightaccent!30},
  decision/.style={diamond,draw=accent,thick,aspect=2.45,align=center,inner sep=1pt,text width=43mm,minimum height=18mm},
  action/.style={rectangle,rounded corners=2mm,draw=darkaccent,align=center,text width=45mm,minimum height=12mm,inner sep=3pt},
  route/.style={->,line width=.75pt,color=darkaccent},
  lab/.style={font=\footnotesize,fill=white,inner sep=1.2pt}
]
\node[start] (start) {Нужно доказать, что\\ четырёхугольник --- параллелограмм};
\node[decision,below=of start] (one) {Есть пара противоположных сторон, равных и параллельных?};
\node[action,right=of one] (yesone) {Признак 1\\ Доказательство завершено};
\node[decision,below=of one] (two) {Обе пары противоположных сторон равны?};
\node[action,right=of two] (yestwo) {Признак 2\\ Доказательство завершено};
\node[decision,below=of two] (three) {Диагонали делятся точкой пересечения пополам?};
\node[action,right=of three] (yesthree) {Признак 3\\ Доказательство завершено};
\node[action,below=of three] (more) {Проведите вспомогательное построение; найдите новые равенства};
\draw[route] (start)--(one);
\draw[route] (one.east)--node[lab,above]{да}(yesone.west);
\draw[route] (one.south)--node[lab,right]{нет}(two.north);
\draw[route] (two.east)--node[lab,above]{да}(yestwo.west);
\draw[route] (two.south)--node[lab,right]{нет}(three.north);
\draw[route] (three.east)--node[lab,above]{да}(yesthree.west);
\draw[route] (three.south)--node[lab,right]{нет}(more.north);
\coordinate (returnleft) at ($(one.west)+(-12mm,0)$);
\draw[route] (more.west) -- (returnleft |- more.west)
  -- node[lab,sloped,above]{проверьте снова} (returnleft |- one.west)
  -- (one.west);
\end{tikzpicture}
\end{center}
\clearpage

\section*{6. Тренировочный блок · Путь А}
\textit{10 заданий. Располагайте вычисления и доказательства последовательно. Для геометрических задач выполняйте собственный чертёж; в заданиях на доказательство обязательно указывайте используемый признак или свойство.}

\begin{enumerate}[label=\textbf{\arabic*.},itemsep=11pt]
\item \textbf{Периметр и отношение.} В параллелограмме \(ABCD\) периметр равен \(84\) см, а \(AB:BC=4:3\). Найдите стороны \(AB\) и \(BC\).
\item \textbf{Углы.} В параллелограмме \(ABCD\) \(\angle A:\angle B=4:5\). Найдите все четыре угла.
\item \textbf{Текстовое условие.} Периметр параллелограмма равен \(70\) см. Сторона \(AB\) на 5 см больше удвоенной стороны \(BC\). Найдите \(AB\) и \(BC\).
\item \textbf{Диагонали: доказательство.} В параллелограмме \(ABCD\) диагонали пересекаются в точке \(O\). На диагонали \(AC\) выбраны точки в порядке \(A,E,O,F,C\) так, что \(AE=CF\). Докажите, что \(BEDF\) --- параллелограмм.
\item \textbf{Средние линии: доказательство.} В произвольном выпуклом четырёхугольнике \(SKLE\) точки \(A,B,C,D\) являются серединами сторон \(SK,KL,LE,ES\) соответственно. Докажите, что \(ABCD\) --- параллелограмм.
\item \textbf{Прямоугольник: доказательство.} В параллелограмме \(ABCD\) точка \(M\) --- середина \(AB\). Известно, что \(MC=MD\). Докажите, что \(ABCD\) --- прямоугольник. \textit{Подсказка:} отметьте середину \(CD\).
\item \textbf{Основание и разряды.} Переведите число \(\sys{3B}{16}\) в десятичную систему и проверьте, что запись допустима.
\item \textbf{Деление с остатком.} Переведите число \(\sys{237}{10}\) в восьмеричную систему. Выполните обратную проверку.
\item \textbf{Тетрады.} Переведите \(\sys{6AF}{16}\) в двоичную систему и определите количество единиц в двоичной записи.
\item \textbf{Триады и тетрады.} Переведите \(\sys{7B4}{16}\) в восьмеричную систему через двоичную. Покажите оба разбиения на группы.
\end{enumerate}

\clearpage
\section*{7. Ответы и ориентиры для самопроверки}
\textit{Сначала завершите собственное решение. Для задач на доказательство сверяйте цепочку рассуждений, а не только конечный вывод.}
\begin{enumerate}[label=\textbf{\arabic*.},itemsep=8pt]
\item \(AB=24\) см, \(BC=18\) см. \(2(4x+3x)=84\Rightarrow x=6\).
\item \(\angle A=\angle C=80^\circ\), \(\angle B=\angle D=100^\circ\).
\item \(BC=10\) см, \(AB=25\) см. \(2[(2x+5)+x]=70\Rightarrow x=10\).
\item \(AO=CO\) и \(AE=CF\Rightarrow EO=FO\). Также \(BO=DO\). Диагонали \(EF\) и \(BD\) делятся пополам в \(O\); признак 3.
\item Проведите диагональ \(SL\). В треугольниках \(SKL\) и \(SEL\) отрезки \(AB\) и \(CD\) --- средние линии: \(AB\parallel SL\), \(CD\parallel SL\), \(AB=CD=SL/2\). Значит, \(AB\parallel CD\) и \(AB=CD\); признак 1.
\item Пусть \(K\) --- середина \(CD\). Из \(AM=DK\), \(AM\parallel DK\) следует, что \(AMKD\) --- параллелограмм, поэтому \(MK\parallel AD\). Так как \(MC=MD\), медиана \(MK\) равнобедренного треугольника \(MCD\) является высотой: \(MK\perp CD\). Значит, \(AD\perp CD\), и \(ABCD\) --- прямоугольник.
\item \(\sys{3B}{16}=3\cdot16+11=\sys{59}{10}\). \(B\) --- допустимая шестнадцатеричная цифра.
\item \(\sys{237}{10}=\sys{355}{8}\). Проверка: \(3\cdot64+5\cdot8+5=237\).
\item \(\sys{6AF}{16}=\sys{0110\ 1010\ 1111}{2}\). Значащих единиц \(2+2+4=8\).
\item \(\sys{7B4}{16}=\sys{0111\ 1011\ 0100}{2}=\sys{011\ 110\ 110\ 100}{2}=\sys{3664}{8}\).
\end{enumerate}
\vfill
\begin{center}\color{darkaccent}\footnotesize Лёвин Артём Александрович · эксклюзивно для Екатерины Скелли\end{center}
\end{document}